\section{OpaqueFunction: lógica Python en el launch}

\subsection{Por qué existe}

Las sustituciones en línea de ROS~2 (como LaunchConfiguration) son mecanismos declarativos: resuelven un valor diferido en el momento del lanzamiento, pero no pueden ejecutar código Python arbitrario. Si el launch file necesita convertir tipos, hacer cálculos, o generar una lista variable de nodos, ese código no cabe en una sustitución.

OpaqueFunction resuelve este límite. Es una acción que, cuando el sistema la procesa, llama a una función Python y le entrega el context\footnote{El context es el objeto que contiene el estado global de la ejecución del launch: los valores actuales de todos los argumentos ya declarados, las variables de entorno, y otras condiciones del sistema. Es el mecanismo por el que la función tiene acceso a los argumentos resueltos.}: el estado global de la ejecución con todos los argumentos ya resueltos. La función retorna una lista de acciones que el sistema agrega a la secuencia de ejecución.

\subsection{El patrón perform(context)}

Dentro de una OpaqueFunction, la forma de leer el valor actual de un argumento es llamar a .perform(context) sobre el objeto LaunchConfiguration correspondiente. Esto fuerza la resolución inmediata del valor y entrega un string, que el código Python puede operar como necesite. El Código~\ref{lst:opaque_patron} muestra el patrón completo:

\begin{figure*}[t]
\centering
\begin{lstlisting}[style=python,
    caption={Patrón OpaqueFunction con perform(context).},
    label={lst:opaque_patron}]
from launch.actions import OpaqueFunction
from launch.substitutions import LaunchConfiguration

def launch_setup(context, *args, **kwargs):
    # .perform(context) extrae el string del argumento
    n = int(LaunchConfiguration('n_robots').perform(context))
    R = float(LaunchConfiguration('neighbor_radius').perform(context))

    actions = []
    for i in range(n):
        actions.append(
            Node(
                package='control_nodes',
                executable='consenso_node',
                namespace=f'robot{i}',
                parameters=[{'neighbor_radius': R}],
            )
        )
    return actions

def generate_launch_description():
    return LaunchDescription([
        DeclareLaunchArgument('n_robots', default_value='3'),
        DeclareLaunchArgument('neighbor_radius', default_value='5.0'),
        OpaqueFunction(function=launch_setup),
    ])
\end{lstlisting}
\end{figure*}


Este patrón es el que usa consenso\_prom.launch.py para generar dinámicamente las acciones de $n$ robots a partir del argumento n\_robots. Sin OpaqueFunction, no sería posible construir un bucle for cuya longitud depende de un argumento de la terminal.

\subsection{Firma de la función}

La firma def launch\_setup(context, *args, **kwargs) es obligatoria. El sistema de launch llama a la función pasando el contexto como primer argumento posicional, y puede pasar argumentos adicionales dependiendo de cómo se configure la acción. El uso de *args, **kwargs garantiza compatibilidad aunque el sistema agregue parámetros en versiones futuras.
